\section{Recolección de datos y gestión del Replay Buffer}

\subsection{Los buenos datos moldean la generalización}

Q-Learning es un método off-policy típico, y durante el entrenamiento se pueden reutilizar transitions históricas, por lo que la calidad de los datos determina directamente el resultado.

\begin{importantbox}{Marco de calidad de datos en tres dimensiones}
\begin{enumerate}
    \item \textbf{Coverage}: si los pares estado-acción cubren toda la distribución de entrenamiento.
    \item \textbf{Freshness}: los datos antiguos pueden quedar obsoletos, sobre todo cuando la política itera muy rápido.
    \item \textbf{Signal-to-Noise}: si la cadena causal entre la recompensa y el siguiente estado es suficientemente clara.
\end{enumerate}
\end{importantbox}

\begin{figure}[H]
\centering
\includegraphics[width=0.86\textwidth,page=1]{06_cs224r_qlearning_2025.pdf}
\caption{Slide: sesgo de distribución en el muestreo del replay buffer y estrategias de corrección}
\end{figure}

\subsection{Repetición de experiencia priorizada y corrección de la distribución}

La repetición de experiencia priorizada (PER) asigna la probabilidad de muestreo según el TD Error, de modo que el entrenamiento se concentre más en las muestras de alto valor o alto error. Los métodos de corrección incluyen:

\begin{itemize}
    \item el importance sampling weight, que se usa para contrarrestar el muestreo no uniforme
    \item los hiperparámetros alpha y beta, que controlan la intensidad del muestreo y de la corrección
\end{itemize}

\begin{knowledgebox}{Consejos de práctica de ingeniería para el Replay Buffer}
1. Limpiar periódicamente los elementos demasiado antiguos para evitar el sesgo de «arranque en frío». 2. Mantener el buffer size en el orden de millones para equilibrar la diversidad de las muestras. 3. Mantener varios buffers para distintas tareas (exploration vs. demonstration).
\end{knowledgebox}

\subsection{Caso: flujo de muestreo de Atari Replay}

Tomando como ejemplo el Q-Learning clásico de Atari, el pipeline de entrenamiento puede simplificarse en los siguientes pasos:

\begin{enumerate}
    \item Recolectar transitions originales mediante epsilon-greedy e insertarlas en el `online buffer`.
    \item Cada 1k steps, generar un mini-batch con prioritized sampling y registrar al mismo tiempo la TD error distribution.
    \item Visualizar el TD error como una curva; si la dirección de la cola larga se desplaza de manera sostenida, activar `replay augmentation` (#32, #64 sampling).
    \item Colocar periódicamente por adelantado los lotes de alta varianza en el `high-priority buffer` para su evaluación y debugging.
\end{enumerate}

\begin{importantbox}{Checklist de monitoreo del Replay}
\begin{itemize}
    \item Registrar el replay hit ratio, y asegurar que la tasa de reutilización del muestreo se mantenga entre 70-95\%.
    \item Cuando el `TD error` aumenta, incrementar automáticamente `beta` para reforzar la corrección del importance sampling.
    \item Copiar las human-verified failure transitions a un buffer separado, para facilitar la inspección posterior.
\end{itemize}
\end{importantbox}

\subsection{Resumen de la sección}

La naturaleza off-policy de Q-Learning hace del replay buffer el motor del rendimiento; la estrategia de muestreo y la calidad de los datos determinan directamente la capacidad de generalización final.

